% HISTORICAL INTEGRATION SNIPPET: promoted on 29 August 2026 after the
% exact-byte production Rethlas pass.

% --------------------------------------------------------------------------
% Verification-boundary paragraph: insert before the existing 29 August
% paragraph for P06.2.
% --------------------------------------------------------------------------
A separate 29 August gate verifies the quasi-smooth scope of P06.1 in
\Cref{thm:p061-quasismooth-ph-rpl}.  Its exact frozen statement and
source-aware proof passed production run
\nolinkurl{20260829T043457Z_0b1b4d161e3b}, item
\nolinkurl{pi_7cd35ebd843cf80a7cd696bb}, using
\texttt{gpt-5.6-sol} at \texttt{xhigh}, final \texttt{correct}, with zero
critical errors and zero gaps.  The checked statement and proof SHA-256
digests are, respectively,
\nolinkurl{0b1b4d161e3b85d342d7be88234a8ed27bfe74be4d24a639ff0456373a7a39e9}
and
\nolinkurl{fdc6abcc4988644656892172fbdbda4178cb86ca710fd0122a90947e598b0b8b}.
The theorem applies to good strict quasi-smooth analytic spaces and requires
no global quasicompactness, separatedness, or boundarylessness.  Together
with the published comparison for arbitrary strict analytic curves
\cite[\S4.1.2 and Proposition~3.15]{GR25Classical}, it leaves P06.1 open only
at non-quasi-smooth points in analytic dimension at least two.
The annular example \Cref{prop:p061-integral-pl-counterexample} shows that
\(\mathbf R\)-PL cannot be replaced by \(\mathbf Q\)-PL or integral PL.

% --------------------------------------------------------------------------
% Provenance record: insert immediately before the existing record
% ``Vanishing CLD current implies harmonicity for real-PL functions''.
% --------------------------------------------------------------------------
\AppLProvenanceRecord
  {Classical pluriharmonic regularity on quasi-smooth analytic spaces}
  {\cref{thm:p061-quasismooth-ph-rpl,prop:p061-integral-pl-counterexample}}
  {Good strict quasi-smooth analytic spaces over complete nontrivially valued
   fields; no global quasicompactness, separatedness, or boundarylessness.
   The proof combines \cite[Proposition~7.8]{GR25Classical},
   \cite[Corollary~3.3.2]{Temkin17}, the rigid \(\acute{e}\)tale sheaf
   criterion \cite[Proposition~3.2.2]{dJvdP96}, and real-PL finite-flat and
   field descent from \cite{GJR26}.  The exact frozen packet passed production
   run \nolinkurl{20260829T043457Z_0b1b4d161e3b}, final \texttt{correct},
   zero critical errors, and zero gaps.  Non-quasi-smooth points in dimension
   at least two remain outside the conclusion.  The \(\sqrt{2}\log|T|\) annular example rules out a
   \(\mathbf Q\)-PL or integral-PL strengthening; no novelty or priority
   claim is made.}

% --------------------------------------------------------------------------
% Questions-not-promoted list: replace the first item by this narrower text.
% --------------------------------------------------------------------------
\item classical pluriharmonic \(\mathbf R\)-PL regularity at non-quasi-smooth points
  in analytic dimension at least two, arbitrary continuous
  higher-dimensional zero-current regularity, and a global CLD--BFJ
  singular-psh equivalence; the quasi-smooth part of the first question is
  closed by \Cref{thm:p061-quasismooth-ph-rpl}, and every strict analytic
  curve is covered by the published curve comparison
  \cite[\S4.1.2 and Proposition~3.15]{GR25Classical}.  The real-PL
  zero-current subproblem is closed by
  \Cref{prop:ch18-real-pl-zero-current-harmonic}, and the proper full-rank
  moment-factor subproblem is closed by
  \Cref{prop:p062-proper-full-rank-moment}, but no theorem supplies that
  factorization for an arbitrary continuous function.  Ordinary CLD-smooth
  harmonic approximation is disproved by
  \Cref{prop:ch19-cld-smooth-harmonic-nonregularization}, while ordinary
  CLD-smooth psh regularization on smooth proper curves is proved by
  \Cref{thm:p063-current-psh-ordinary-regularization};
